\section{Beurteilung des Kursverlaufs}
\label{sec:kurs-rdy}

\subsection{Was die Reihe zeigt, und was an ihr bereinigt ist}

Auch hier zuerst die Bereinigung. Zum 28.~Oktober 2024 wurde die Aktie im
Verh\"altnis eins zu f\"unf geteilt, das Verh\"altnis des
Hinterlegungsscheins zur Aktie blieb dabei
unver\"andert.\quelleZF{2026}{157}\quelleMerken{s31zfxcacgxbfh}
Der Schein kostete am Tag vor dem Split also das F\"unffache dessen, was er am
Tag danach kostete. Die hier verwendete Reihe ist zur\"uckgerechnet; um den
Splittermin herum findet sich in Abbildung~\ref{fig:kurs-rdy} kein Sprung, und
das ist der Beleg daf\"ur, dass die Bereinigung greift.

\"Uber die zehn Jahre bis zum \Datum{\RdyKursTag} stieg der Kurs von
\USDje{\RdyKursSchlussXVI} auf \USDje{\RdyKurs}. Das sind
\Pct{\RdyRenditeZehn} pro Jahr, gegen eine H\"urde von \Pct{\Huerde}. Die
beiden Gr\"o{\ss}en sind nicht deckungsgleich: Die H\"urde ist eine
Gesamtrendite, die Kursreihe enth\"alt keine Dividenden. Die
Dividendenrendite dieses Titels liegt bei \Pct{\RdyDividendenrendite}, und
damit bleibt der Abstand bestehen.

\bk{Der Titel hat den Branchenindex \"uber zehn Jahre also nicht geschlagen,
sondern deutlich verfehlt, und zwar \"uber einen Zeitraum, der lang genug
ist, um nicht als Momentaufnahme zu gelten.}

\subsection{Was zu erkl\"aren ist: der R\"uckgang seit Juni 2026}

Am \Datum{\RdyKurstagHoch} stand der Schein bei \USDje{\RdyKursHochXXVI}; heute
steht er \PctBetrag{\RdyAbstandHoch} darunter und
\PctBetrag{\RdyRenditeLaufend} unter dem Jahresschluss 2025.

\subsection{Was gefallen ist: Ertrag, nicht Bewertung}

\bk{Anders als bei AstraZeneca ist die Erkl\"arung hier nicht die
Bewertung, sondern der Ertrag, und das ist der ung\"unstigere der beiden
Befunde.}

\bk{Der Reihe nach. Im Gesch\"aftsjahr FY2026 fiel der Jahres\"uberschuss um
\PctBetrag{\RdyErgebnisWachstum}, obwohl der Umsatz stieg; die
Rohertragsmarge sank um mehr als f\"unf Prozentpunkte
(Abschnitt~\ref{sec:abschluss-rdy}). Im Quartal danach fiel der Gewinn um
\PctBetrag{\RdyQuartalGewinnWachstum} und der operative Mittelzufluss wurde
negativ (Abschnitt~\ref{sec:strategie-rdy}). Der Kurs ist diesen Zahlen
gefolgt, nicht vorausgelaufen.}

\bk{Die Ursache benennt das Unternehmen selbst, wenn auch nicht in einer
Summe: Preiserosion bestehender Produkte in Nordamerika einschlie{\ss}lich
Lenalidomid. Der nordamerikanische Umsatz fiel um
\MioINR{\RdyNordamerikaRueckgang}, das Therapiegebiet Onkologie um
\MioINR{\RdyOnkologieRueckgang}. Was an die Stelle getreten ist (das
zugekaufte Nikotinersatzgesch\"aft mit \MioINR{\RdyNrtUmsatz}), ersetzt
einen Teil des Umsatzes und ist ein Konsumg\"utergesch\"aft mit anderer
Marge; der Rohertrag des Segments fiel im selben Jahr um
\Pct{\RdyRohertragGgWachstum}.}

\bk{Damit ist die Trennung, die Teil~\ref{sec:kurs-azn} f\"ur AstraZeneca
vornehmen musste, hier einfach: Es ist kein Bewertungsereignis. Das
Kurs-Gewinn-Verh\"altnis lag am Jahreshoch bei \Faktor{\RdyKgvHoch} und liegt
heute bei \Faktor{\RdyKursGewinn}; es ist kaum gesunken, weil der Nenner
mitgefallen ist. Der Markt zahlt f\"ur dieses Unternehmen dasselbe Vielfache
wie im Juni; es verdient nur weniger.}

\subsection{Was der Kurs nicht abbildet}

\bk{Was der Kurs nicht abbilden kann, weil es niemand ver\"offentlicht hat:
wie viel von dem, was gefallen ist, einmalig war und wie viel bleibt. Der
Regalbestandsausgleich \"uber \MioINR{\RdySsaEffekt} ist ein Einmaleffekt und
als solcher beziffert. Die Preiserosion ist es nicht, und der Umsatz, der am
auslaufenden Lenalidomid-Gesch\"aft h\"angt, ist \"uberhaupt nicht beziffert
(Abschnitt~\ref{sec:lenalidomid-rdy}). Ein Leser, der wissen will, ob
\INRje{\RdyErgebnisJeAktie} je Aktie der neue Boden oder eine Zwischenstufe
sind, findet in den Prim\"arquellen keine Antwort.}
